% Hurewicz theorem: pi_n -> H_n via fundamental class.
\begin{tikzpicture}[
  every node/.style={font=\small},
  arrow/.style={-{Stealth[length=2.5mm]}, thick},
]
  % Sphere on the left
  \draw[blue!60!black, thick, fill=blue!8] (-4.5, 0) circle (1.0cm);
  \node[blue!60!black] at (-4.5, 0) {$S^n$};
  \node[font=\footnotesize, gray!50!black] at (-4.5, -1.4) {with $[S^n]$ generator};

  % Map sigma into X
  \draw[arrow] (-3.4, 0) -- node[above] {$\sigma$} (-1.6, 0);

  % Space X
  \draw[gray!50!black, very thick, fill=gray!8] plot[smooth cycle]
    coordinates {(-1.5, 0.6) (-0.0, 1.0) (1.5, 0.5) (1.7, -0.4) (0.3, -1.2) (-1.2, -0.6)};
  \node at (1.0, 0.6) {$X$};

  % Image of sphere inside X (a wiggly sphere outline)
  \draw[blue!75!black, thick] (0.0, -0.2) circle (0.5cm);

  % Equation showing the Hurewicz map
  \node[anchor=west] at (3.5, 0.8) {Hurewicz map:};
  \node[anchor=west, font=\small] at (3.5, 0.0) {$h_n : \pi_n(X, x_0) \;\to\; H_n(X)$};
  \node[anchor=west, font=\small] at (3.5, -0.8) {$[\sigma] \;\mapsto\; \sigma_*[S^n]$};

  % Theorem text
  \node[align=left, anchor=west, font=\footnotesize] at (-5.0, -2.4)
    {{\bfseries Hurewicz theorem.} If $X$ is $(n-1)$-connected ($n \geq 2$), then\\
     $\quad H_k(X) = 0$ for $1 \leq k < n$, \quad and \quad $h_n : \pi_n(X) \xrightarrow{\;\sim\;} H_n(X)$.\\
     For $n = 1$: $h_1$ is the abelianisation map $\pi_1(X) \to H_1(X) = \pi_1(X)^{\mathrm{ab}}$.};

  \begin{scope}[on background layer]
    \fill[teal!4, rounded corners=12pt] (-6.0, -3.6) rectangle (10.5, 1.8);
  \end{scope}
\end{tikzpicture}
